%!TeX root=../persuasiontop.tex
\chapter[Chapter \thechapter]{}
	
	\begin{figure}[t!]
\centering
\includegraphics[width=\linewidth]{nouveau12}
\end{figure}

\lettrine[lines=4,lraise=0.3]{A}{nne} and Henrietta, finding themselves the earliest of the party the next morning, agreed to stroll down to the sea before breakfast. They went to the sands, to watch the flowing of the tide, which a fine south-easterly breeze was bringing in with all the grandeur which so flat a shore admitted. They praised the morning; gloried in the sea; sympathized in the delight of the fresh-feeling breeze—and were silent; till Henrietta suddenly began again with—

»Oh! yes,—I am quite convinced that, with very few exceptions, the sea-air always does good. There can be no doubt of its having been of the greatest service to Dr Shirley, after his illness, last spring twelvemonth. He declares himself, that coming to Lyme for a month, did him more good than all the medicine he took; and that being by the sea always makes him feel young again. Now, I cannot help thinking it a pity that he does not live entirely by the sea. I do think he had better leave Uppercross entirely, and fix at Lyme. Do not you, Anne? Do not you agree with me, that it is the best thing he could do, both for himself and Mrs Shirley? She has cousins here, you know, and many acquaintance, which would make it cheerful for her, and I am sure she would be glad to get to a place where she could have medical attendance at hand, in case of his having another seizure. Indeed I think it quite melancholy to have such excellent people as Dr and Mrs Shirley, who have been doing good all their lives, wearing out their last days in a place like Uppercross, where, excepting our family, they seem shut out from all the world. I wish his friends would propose it to him. I really think they ought. And, as to procuring a dispensation, there could be no difficulty at his time of life, and with his character. My only doubt is, whether anything could persuade him to leave his parish. He is so very strict and scrupulous in his notions; over-scrupulous I must say. Do not you think, Anne, it is being over-scrupulous? Do not you think it is quite a mistaken point of conscience, when a clergyman sacrifices his health for the sake of duties, which may be just as well performed by another person? And at Lyme too, only seventeen miles off, he would be near enough to hear, if people thought there was anything to complain of.«

Anne smiled more than once to herself during this speech, and entered into the subject, as ready to do good by entering into the feelings of a young lady as of a young man, though here it was good of a lower standard, for what could be offered but general acquiescence? She said all that was reasonable and proper on the business; felt the claims of Dr Shirley to repose as she ought; saw how very desirable it was that he should have some active, respectable young man as a resident curate, and was even courteous enough to hint at the advantage of such resident curate's being married.

»I wish,« said Henrietta, very well pleased with her companion, »I wish Lady Russell lived at Uppercross, and were intimate with Dr Shirley. I have always heard of Lady Russell as a woman of the greatest influence with everybody! I always look upon her as able to persuade a person to anything! I am afraid of her, as I have told you before, quite afraid of her, because she is so very clever; but I respect her amazingly, and wish we had such a neighbour at Uppercross.«

Anne was amused by Henrietta's manner of being grateful, and amused also that the course of events and the new interests of Henrietta's views should have placed her friend at all in favour with any of the Musgrove family; she had only time, however, for a general answer, and a wish that such another woman were at Uppercross, before all subjects suddenly ceased, on seeing Louisa and Captain Wentworth coming towards them. They came also for a stroll till breakfast was likely to be ready; but Louisa recollecting immediately afterwards that she had something to procure at a shop, invited them all to go back with her into the town. They were all at her disposal.

When they came to the steps, leading upwards from the beach, a gentleman, at the same moment preparing to come down, politely drew back, and stopped to give them way. They ascended and passed him; and as they passed, Anne's face caught his eye, and he looked at her with a degree of earnest admiration, which she could not be insensible of. She was looking remarkably well; her very regular, very pretty features, having the bloom and freshness of youth restored by the fine wind which had been blowing on her complexion, and by the animation of eye which it had also produced. It was evident that the gentleman, (completely a gentleman in manner) admired her exceedingly. Captain Wentworth looked round at her instantly in a way which shewed his noticing of it. He gave her a momentary glance, a glance of brightness, which seemed to say, »That man is struck with you, and even I, at this moment, see something like Anne Elliot again.«

After attending Louisa through her business, and loitering about a little longer, they returned to the inn; and Anne, in passing afterwards quickly from her own chamber to their dining-room, had nearly run against the very same gentleman, as he came out of an adjoining apartment. She had before conjectured him to be a stranger like themselves, and determined that a well-looking groom, who was strolling about near the two inns as they came back, should be his servant. Both master and man being in mourning assisted the idea. It was now proved that he belonged to the same inn as themselves; and this second meeting, short as it was, also proved again by the gentleman's looks, that he thought hers very lovely, and by the readiness and propriety of his apologies, that he was a man of exceedingly good manners. He seemed about thirty, and though not handsome, had an agreeable person. Anne felt that she should like to know who he was.

They had nearly done breakfast, when the sound of a carriage, (almost the first they had heard since entering Lyme) drew half the party to the window. It was a gentleman's carriage, a curricle, but only coming round from the stable-yard to the front door; somebody must be going away. It was driven by a servant in mourning.

The word curricle made Charles Musgrove jump up that he might compare it with his own; the servant in mourning roused Anne's curiosity, and the whole six were collected to look, by the time the owner of the curricle was to be seen issuing from the door amidst the bows and civilities of the household, and taking his seat, to drive off.

»Ah!« cried Captain Wentworth, instantly, and with half a glance at Anne, »it is the very man we passed.«

The Miss Musgroves agreed to it; and having all kindly watched him as far up the hill as they could, they returned to the breakfast table. The waiter came into the room soon afterwards.

»Pray,« said Captain Wentworth, immediately, »can you tell us the name of the gentleman who is just gone away?«

»Yes, Sir, a Mr Elliot, a gentleman of large fortune, came in last night from Sidmouth. Dare say you heard the carriage, sir, while you were at dinner; and going on now for Crewkherne, in his way to Bath and London.«

»Elliot!« Many had looked on each other, and many had repeated the name, before all this had been got through, even by the smart rapidity of a waiter.

»Bless me!« cried Mary; »it must be our cousin; it must be our Mr Elliot, it must, indeed! Charles, Anne, must not it? In mourning, you see, just as our Mr Elliot must be. How very extraordinary! In the very same inn with us! Anne, must not it be our Mr Elliot? my father's next heir? Pray sir,« turning to the waiter, »did not you hear, did not his servant say whether he belonged to the Kellynch family?«

»No, ma'am, he did not mention no particular family; but he said his master was a very rich gentleman, and would be a baronight some day.«

»There! you see!« cried Mary in an ecstasy, »just as I said! Heir to Sir Walter Elliot! I was sure that would come out, if it was so. Depend upon it, that is a circumstance which his servants take care to publish, wherever he goes. But, Anne, only conceive how extraordinary! I wish I had looked at him more. I wish we had been aware in time, who it was, that he might have been introduced to us. What a pity that we should not have been introduced to each other! Do you think he had the Elliot countenance? I hardly looked at him, I was looking at the horses; but I think he had something of the Elliot countenance, I wonder the arms did not strike me! Oh! the great-coat was hanging over the panel, and hid the arms, so it did; otherwise, I am sure, I should have observed them, and the livery too; if the servant had not been in mourning, one should have known him by the livery.«

»Putting all these very extraordinary circumstances together,« said Captain Wentworth, »we must consider it to be the arrangement of Providence, that you should not be introduced to your cousin.«

When she could command Mary's attention, Anne quietly tried to convince her that their father and Mr Elliot had not, for many years, been on such terms as to make the power of attempting an introduction at all desirable.

At the same time, however, it was a secret gratification to herself to have seen her cousin, and to know that the future owner of Kellynch was undoubtedly a gentleman, and had an air of good sense. She would not, upon any account, mention her having met with him the second time; luckily Mary did not much attend to their having passed close by him in their earlier walk, but she would have felt quite ill-used by Anne's having actually run against him in the passage, and received his very polite excuses, while she had never been near him at all; no, that cousinly little interview must remain a perfect secret.

»Of course,« said Mary, »you will mention our seeing Mr Elliot, the next time you write to Bath. I think my father certainly ought to hear of it; do mention all about him.«

Anne avoided a direct reply, but it was just the circumstance which she considered as not merely unnecessary to be communicated, but as what ought to be suppressed. The offence which had been given her father, many years back, she knew; Elizabeth's particular share in it she suspected; and that Mr Elliot's idea always produced irritation in both was beyond a doubt. Mary never wrote to Bath herself; all the toil of keeping up a slow and unsatisfactory correspondence with Elizabeth fell on Anne.

Breakfast had not been long over, when they were joined by Captain and Mrs Harville and Captain Benwick; with whom they had appointed to take their last walk about Lyme. They ought to be setting off for Uppercross by one, and in the meanwhile were to be all together, and out of doors as long as they could.

Anne found Captain Benwick getting near her, as soon as they were all fairly in the street. Their conversation the preceding evening did not disincline him to seek her again; and they walked together some time, talking as before of Mr Scott and Lord Byron, and still as unable as before, and as unable as any other two readers, to think exactly alike of the merits of either, till something occasioned an almost general change amongst their party, and instead of Captain Benwick, she had Captain Harville by her side.

»Miss Elliot,« said he, speaking rather low, »you have done a good deed in making that poor fellow talk so much. I wish he could have such company oftener. It is bad for him, I know, to be shut up as he is; but what can we do? We cannot part.«

»No,« said Anne, »that I can easily believe to be impossible; but in time, perhaps—we know what time does in every case of affliction, and you must remember, Captain Harville, that your friend may yet be called a young mourner—only last summer, I understand.«

»Ay, true enough,« (with a deep sigh) »only June.«

»And not known to him, perhaps, so soon.«

»Not till the first week of August, when he came home from the Cape, just made into the Grappler. I was at Plymouth dreading to hear of him; he sent in letters, but the Grappler was under orders for Portsmouth. There the news must follow him, but who was to tell it? not I. I would as soon have been run up to the yard-arm. Nobody could do it, but that good fellow« (pointing to Captain Wentworth). »The Laconia had come into Plymouth the week before; no danger of her being sent to sea again. He stood his chance for the rest; wrote up for leave of absence, but without waiting the return, travelled night and day till he got to Portsmouth, rowed off to the Grappler that instant, and never left the poor fellow for a week. That's what he did, and nobody else could have saved poor James. You may think, Miss Elliot, whether he is dear to us!«

Anne did think on the question with perfect decision, and said as much in reply as her own feeling could accomplish, or as his seemed able to bear, for he was too much affected to renew the subject, and when he spoke again, it was of something totally different.

Mrs Harville's giving it as her opinion that her husband would have quite walking enough by the time he reached home, determined the direction of all the party in what was to be their last walk; they would accompany them to their door, and then return and set off themselves. By all their calculations there was just time for this; but as they drew near the Cobb, there was such a general wish to walk along it once more, all were so inclined, and Louisa soon grew so determined, that the difference of a quarter of an hour, it was found, would be no difference at all; so with all the kind leave-taking, and all the kind interchange of invitations and promises which may be imagined, they parted from Captain and Mrs Harville at their own door, and still accompanied by Captain Benwick, who seemed to cling to them to the last, proceeded to make the proper adieus to the Cobb.

Anne found Captain Benwick again drawing near her. Lord Byron's »dark blue seas« could not fail of being brought forward by their present view, and she gladly gave him all her attention as long as attention was possible. It was soon drawn, perforce another way.

There was too much wind to make the high part of the new Cobb pleasant for the ladies, and they agreed to get down the steps to the lower, and all were contented to pass quietly and carefully down the steep flight, excepting Louisa; she must be jumped down them by Captain Wentworth. In all their walks, he had had to jump her from the stiles; the sensation was delightful to her. The hardness of the pavement for her feet, made him less willing upon the present occasion; he did it, however. She was safely down, and instantly, to show her enjoyment, ran up the steps to be jumped down again. He advised her against it, thought the jar too great; but no, he reasoned and talked in vain, she smiled and said, »I am determined I will:« he put out his hands; she was too precipitate by half a second, she fell on the pavement on the Lower Cobb, and was taken up lifeless! There was no wound, no blood, no visible bruise; but her eyes were closed, she breathed not, her face was like death. The horror of the moment to all who stood around!

Captain Wentworth, who had caught her up, knelt with her in his arms, looking on her with a face as pallid as her own, in an agony of silence. »She is dead! she is dead!« screamed Mary, catching hold of her husband, and contributing with his own horror to make him immoveable; and in another moment, Henrietta, sinking under the conviction, lost her senses too, and would have fallen on the steps, but for Captain Benwick and Anne, who caught and supported her between them.

»Is there no one to help me?« were the first words which burst from Captain Wentworth, in a tone of despair, and as if all his own strength were gone.

»Go to him, go to him,« cried Anne, »for heaven's sake go to him. I can support her myself. Leave me, and go to him. Rub her hands, rub her temples; here are salts; take them, take them.«

Captain Benwick obeyed, and Charles at the same moment, disengaging himself from his wife, they were both with him; and Louisa was raised up and supported more firmly between them, and everything was done that Anne had prompted, but in vain; while Captain Wentworth, staggering against the wall for his support, exclaimed in the bitterest agony—

»Oh God! her father and mother!«

»A surgeon!« said Anne.

He caught the word; it seemed to rouse him at once, and saying only—»True, true, a surgeon this instant,« was darting away, when Anne eagerly suggested—

»Captain Benwick, would not it be better for Captain Benwick? He knows where a surgeon is to be found.«

Every one capable of thinking felt the advantage of the idea, and in a moment (it was all done in rapid moments) Captain Benwick had resigned the poor corpse-like figure entirely to the brother's care, and was off for the town with the utmost rapidity.

As to the wretched party left behind, it could scarcely be said which of the three, who were completely rational, was suffering most: Captain Wentworth, Anne, or Charles, who, really a very affectionate brother, hung over Louisa with sobs of grief, and could only turn his eyes from one sister, to see the other in a state as insensible, or to witness the hysterical agitations of his wife, calling on him for help which he could not give.

Anne, attending with all the strength and zeal, and thought, which instinct supplied, to Henrietta, still tried, at intervals, to suggest comfort to the others, tried to quiet Mary, to animate Charles, to assuage the feelings of Captain Wentworth. Both seemed to look to her for directions.

»Anne, Anne,« cried Charles, »What is to be done next? What, in heaven's name, is to be done next?«

Captain Wentworth's eyes were also turned towards her.

»Had not she better be carried to the inn? Yes, I am sure: carry her gently to the inn.«

»Yes, yes, to the inn,« repeated Captain Wentworth, comparatively collected, and eager to be doing something. »I will carry her myself. Musgrove, take care of the others.«

By this time the report of the accident had spread among the workmen and boatmen about the Cobb, and many were collected near them, to be useful if wanted, at any rate, to enjoy the sight of a dead young lady, nay, two dead young ladies, for it proved twice as fine as the first report. To some of the best-looking of these good people Henrietta was consigned, for, though partially revived, she was quite helpless; and in this manner, Anne walking by her side, and Charles attending to his wife, they set forward, treading back with feelings unutterable, the ground, which so lately, so very lately, and so light of heart, they had passed along.

They were not off the Cobb, before the Harvilles met them. Captain Benwick had been seen flying by their house, with a countenance which showed something to be wrong; and they had set off immediately, informed and directed as they passed, towards the spot. Shocked as Captain Harville was, he brought senses and nerves that could be instantly useful; and a look between him and his wife decided what was to be done. She must be taken to their house; all must go to their house; and await the surgeon's arrival there. They would not listen to scruples: he was obeyed; they were all beneath his roof; and while Louisa, under Mrs Harville's direction, was conveyed up stairs, and given possession of her own bed, assistance, cordials, restoratives were supplied by her husband to all who needed them.

Louisa had once opened her eyes, but soon closed them again, without apparent consciousness. This had been a proof of life, however, of service to her sister; and Henrietta, though perfectly incapable of being in the same room with Louisa, was kept, by the agitation of hope and fear, from a return of her own insensibility. Mary, too, was growing calmer.

The surgeon was with them almost before it had seemed possible. They were sick with horror, while he examined; but he was not hopeless. The head had received a severe contusion, but he had seen greater injuries recovered from: he was by no means hopeless; he spoke cheerfully.

That he did not regard it as a desperate case, that he did not say a few hours must end it, was at first felt, beyond the hope of most; and the ecstasy of such a reprieve, the rejoicing, deep and silent, after a few fervent ejaculations of gratitude to Heaven had been offered, may be conceived.

The tone, the look, with which »Thank God!« was uttered by Captain Wentworth, Anne was sure could never be forgotten by her; nor the sight of him afterwards, as he sat near a table, leaning over it with folded arms and face concealed, as if overpowered by the various feelings of his soul, and trying by prayer and reflection to calm them.

Louisa's limbs had escaped. There was no injury but to the head.

It now became necessary for the party to consider what was best to be done, as to their general situation. They were now able to speak to each other and consult. That Louisa must remain where she was, however distressing to her friends to be involving the Harvilles in such trouble, did not admit a doubt. Her removal was impossible. The Harvilles silenced all scruples; and, as much as they could, all gratitude. They had looked forward and arranged everything before the others began to reflect. Captain Benwick must give up his room to them, and get another bed elsewhere; and the whole was settled. They were only concerned that the house could accommodate no more; and yet perhaps, by »putting the children away in the maid's room, or swinging a cot somewhere,« they could hardly bear to think of not finding room for two or three besides, supposing they might wish to stay; though, with regard to any attendance on Miss Musgrove, there need not be the least uneasiness in leaving her to Mrs Harville's care entirely. Mrs Harville was a very experienced nurse, and her nursery-maid, who had lived with her long, and gone about with her everywhere, was just such another. Between these two, she could want no possible attendance by day or night. And all this was said with a truth and sincerity of feeling irresistible.

Charles, Henrietta, and Captain Wentworth were the three in consultation, and for a little while it was only an interchange of perplexity and terror. »Uppercross, the necessity of some one's going to Uppercross; the news to be conveyed; how it could be broken to Mr and Mrs Musgrove; the lateness of the morning; an hour already gone since they ought to have been off; the impossibility of being in tolerable time.« At first, they were capable of nothing more to the purpose than such exclamations; but, after a while, Captain Wentworth, exerting himself, said—

»We must be decided, and without the loss of another minute. Every minute is valuable. Some one must resolve on being off for Uppercross instantly. Musgrove, either you or I must go.«

Charles agreed, but declared his resolution of not going away. He would be as little incumbrance as possible to Captain and Mrs Harville; but as to leaving his sister in such a state, he neither ought, nor would. So far it was decided; and Henrietta at first declared the same. She, however, was soon persuaded to think differently. The usefulness of her staying! She who had not been able to remain in Louisa's room, or to look at her, without sufferings which made her worse than helpless! She was forced to acknowledge that she could do no good, yet was still unwilling to be away, till, touched by the thought of her father and mother, she gave it up; she consented, she was anxious to be at home.

The plan had reached this point, when Anne, coming quietly down from Louisa's room, could not but hear what followed, for the parlour door was open.

»Then it is settled, Musgrove,« cried Captain Wentworth, »that you stay, and that I take care of your sister home. But as to the rest, as to the others, if one stays to assist Mrs Harville, I think it need be only one. Mrs Charles Musgrove will, of course, wish to get back to her children; but if Anne will stay, no one so proper, so capable as Anne.«

She paused a moment to recover from the emotion of hearing herself so spoken of. The other two warmly agreed with what he said, and she then appeared.

»You will stay, I am sure; you will stay and nurse her;« cried he, turning to her and speaking with a glow, and yet a gentleness, which seemed almost restoring the past. She coloured deeply, and he recollected himself and moved away. She expressed herself most willing, ready, happy to remain. »It was what she had been thinking of, and wishing to be allowed to do. A bed on the floor in Louisa's room would be sufficient for her, if Mrs Harville would but think so.«

One thing more, and all seemed arranged. Though it was rather desirable that Mr and Mrs Musgrove should be previously alarmed by some share of delay; yet the time required by the Uppercross horses to take them back, would be a dreadful extension of suspense; and Captain Wentworth proposed, and Charles Musgrove agreed, that it would be much better for him to take a chaise from the inn, and leave Mr Musgrove's carriage and horses to be sent home the next morning early, when there would be the farther advantage of sending an account of Louisa's night.

Captain Wentworth now hurried off to get everything ready on his part, and to be soon followed by the two ladies. When the plan was made known to Mary, however, there was an end of all peace in it. She was so wretched and so vehement, complained so much of injustice in being expected to go away instead of Anne; Anne, who was nothing to Louisa, while she was her sister, and had the best right to stay in Henrietta's stead! Why was not she to be as useful as Anne? And to go home without Charles, too, without her husband! No, it was too unkind. And in short, she said more than her husband could long withstand, and as none of the others could oppose when he gave way, there was no help for it; the change of Mary for Anne was inevitable.

Anne had never submitted more reluctantly to the jealous and ill-judging claims of Mary; but so it must be, and they set off for the town, Charles taking care of his sister, and Captain Benwick attending to her. She gave a moment's recollection, as they hurried along, to the little circumstances which the same spots had witnessed earlier in the morning. There she had listened to Henrietta's schemes for Dr Shirley's leaving Uppercross; farther on, she had first seen Mr Elliot; a moment seemed all that could now be given to any one but Louisa, or those who were wrapped up in her welfare.

Captain Benwick was most considerately attentive to her; and, united as they all seemed by the distress of the day, she felt an increasing degree of good-will towards him, and a pleasure even in thinking that it might, perhaps, be the occasion of continuing their acquaintance.

Captain Wentworth was on the watch for them, and a chaise and four in waiting, stationed for their convenience in the lowest part of the street; but his evident surprise and vexation at the substitution of one sister for the other, the change in his countenance, the astonishment, the expressions begun and suppressed, with which Charles was listened to, made but a mortifying reception of Anne; or must at least convince her that she was valued only as she could be useful to Louisa.

She endeavoured to be composed, and to be just. Without emulating the feelings of an Emma towards her Henry, she would have attended on Louisa with a zeal above the common claims of regard, for his sake; and she hoped he would not long be so unjust as to suppose she would shrink unnecessarily from the office of a friend.

In the meanwhile she was in the carriage. He had handed them both in, and placed himself between them; and in this manner, under these circumstances, full of astonishment and emotion to Anne, she quitted Lyme. How the long stage would pass; how it was to affect their manners; what was to be their sort of intercourse, she could not foresee. It was all quite natural, however. He was devoted to Henrietta; always turning towards her; and when he spoke at all, always with the view of supporting her hopes and raising her spirits. In general, his voice and manner were studiously calm. To spare Henrietta from agitation seemed the governing principle. Once only, when she had been grieving over the last ill-judged, ill-fated walk to the Cobb, bitterly lamenting that it ever had been thought of, he burst forth, as if wholly overcome—

»Don't talk of it, don't talk of it,« he cried. »Oh God! that I had not given way to her at the fatal moment! Had I done as I ought! But so eager and so resolute! Dear, sweet Louisa!«

Anne wondered whether it ever occurred to him now, to question the justness of his own previous opinion as to the universal felicity and advantage of firmness of character; and whether it might not strike him that, like all other qualities of the mind, it should have its proportions and limits. She thought it could scarcely escape him to feel that a persuadable temper might sometimes be as much in favour of happiness as a very resolute character.

They got on fast. Anne was astonished to recognise the same hills and the same objects so soon. Their actual speed, heightened by some dread of the conclusion, made the road appear but half as long as on the day before. It was growing quite dusk, however, before they were in the neighbourhood of Uppercross, and there had been total silence among them for some time, Henrietta leaning back in the corner, with a shawl over her face, giving the hope of her having cried herself to sleep; when, as they were going up their last hill, Anne found herself all at once addressed by Captain Wentworth. In a low, cautious voice, he said:—

»I have been considering what we had best do. She must not appear at first. She could not stand it. I have been thinking whether you had not better remain in the carriage with her, while I go in and break it to Mr and Mrs Musgrove. Do you think this is a good plan?«

She did: he was satisfied, and said no more. But the remembrance of the appeal remained a pleasure to her, as a proof of friendship, and of deference for her judgement, a great pleasure; and when it became a sort of parting proof, its value did not lessen.

When the distressing communication at Uppercross was over, and he had seen the father and mother quite as composed as could be hoped, and the daughter all the better for being with them, he announced his intention of returning in the same carriage to Lyme; and when the horses were baited, he was off.